\section{Introduction}\label{chap:intro}

\subsection{Background}

Daniell-style integration theory \cite{Daniell1918} does not construct a measure
first, but starts from a positive continuous linear functional on a Riesz space
of functions.
From this standpoint an integral is strictly an operation assigning real numbers
to functions on a base space, and it requires no additional structure, such as a
$\sigma$-algebra, on the base space.
The advantage of this standpoint is that it separates out how much of what is
needed to construct the Lebesgue integral belongs to the side of the base space
and where the side of the function space and the functional begins.

This advantage is, however, usually discussed as an argument internal to a
single integration system.
How integration systems on different base spaces are to be related to one
another, and how the Daniell extension behaves with respect to that relation,
are not normally made an explicit subject.
The purpose of this paper is to organize that part in the language of
categories.

Note that requiring no structure on the base space was also the historical
motivation for this standpoint.
Daniell himself applied this construction to infinite-dimensional base spaces
\cite{Daniell1919}.

\subsection{Organization and contributions of this paper}

In this paper the examples are not collected into an independent section but
placed immediately after the definition each of them tests.
Examples of integration systems are
Example~\ref{ex:counting} and Example~\ref{ex:step-functions}; that the Daniell
extension recovers the Lebesgue integral is
Example~\ref{ex:lebesgue}; that operator morphisms include stochastic matrices
is Example~\ref{ex:stochastic-matrix}.
In Subsection~\ref{sec:measure} we also confirm, via the Stone condition and the
Daniell--Stone representation theorem, that integration systems recover measure
theory.
The results of this paper do not use that condition.
That is, this paper \textbf{does not assume the Stone condition and does not
take representation by a measure as a premise of the theory}.

Section~\ref{chap:preliminaries} is preliminary.
It defines integration systems $(X,L,I)$, integrability via series
representations, the Daniell extension $(X,L^1,I^1)$, and the monotone
convergence theorem and Daniell closedness, and lists the statements needed
later.
\textbf{Nothing in this section is new, and as a rule we refer to earlier work
for the proofs as well.}

Section~\ref{chap:functoriality} is the main subject of this paper.
We introduce the category $\mathbf{Int}$ of integration systems and their
morphisms, the full subcategory $\mathbf{MInt}$ of integration systems
satisfying the monotone convergence theorem, and the further full subcategory
$\mathbf{DInt}$ of Daniell closed ones, and show that the Daniell extension
defines a functor
\[
  \mathcal{D}:\mathbf{Int}\to\mathbf{DInt}
\]
and that this is left adjoint to the inclusion functor $\iota_{D}$.
That is, $\mathbf{DInt}$ is a reflective subcategory of $\mathbf{Int}$ and the
Daniell extension is its reflector.

We then show, by a counterexample on a finite set, that the restriction imposing
Daniell closedness on the codomain in this theorem cannot be dispensed with.
That is, the functor $\iota_M\circ\mathcal{D}:\mathbf{Int}\to\mathbf{MInt}$
defined by the Daniell extension is \emph{not} left adjoint to the inclusion
functor $\iota:\mathbf{MInt}\to\mathbf{Int}$.
We then show that $\iota$ does have a left adjoint, but that its reflector is
\emph{not} the Daniell extension
(Theorem~\ref{thm:mint-inclusion-has-left-adjoint}).
We investigate when these two reflections agree.
Using the $\sigma$-algebra $\sigma(L)$ generated by $L$, the second reflection
$L^{\mathrm{M}}$ is contained in the $\sigma(L)$-measurable part of $L^1$
(Proposition~\ref{prop:LM-in-Lsigma}), and under the Stone condition the two
coincide (Theorem~\ref{thm:LM-equals-Lsigma}).
This characterizes the agreement of the two reflections as a property of
$(X,L,I)$ itself (Corollary~\ref{cor:when-coincide-sigma}).
As a corollary, for the integration system consisting of the step functions on
$[0,1)$ and the Riemann integral, $L^{\mathrm{M}}$ is the set of all Borel
measurable integrable functions and $L^1$ is the set of all Lebesgue integrable
functions, so the two reflections do not agree
(Corollary~\ref{cor:lebesgue-reflections-differ}).
Furthermore, for an arbitrary integration system, $L^1$ splits as the sum of its
$\sigma(L)$-measurable part and the part supported on null sets
(Proposition~\ref{prop:l1-decomposition}).
Hence \textbf{the gap between the two reflections amounts exactly to the
functions supported on null sets}, and in the case of step functions this
appears as the difference between Borel and Lebesgue, that is, as the completion
of the measure.
Moreover, \textbf{neither separation of points by $L$ nor the Stone condition is
a sufficient condition for the two reflections to agree}.
The remaining question is whether the Stone condition is necessary
(Problem~\ref{prob:when-reflections-coincide}).
Finally we prove a main theorem of the same shape in the category obtained by
dividing morphisms by null functions (Theorem~\ref{thm:main-ae}).

In Section~\ref{chap:operator} we detach morphisms from maps of base spaces and
redefine them as positive linear maps between Riesz spaces.
This is based on the operators introduced in \cite{Yamagami} in the context of
iterated integrals, but with the order continuity imposed there removed.
The reason for removing it is given at the beginning of
Section~\ref{chap:operator} and in
Remark~\ref{rem:order-continuity-redundant}.
We show that in this category an operator morphism extends uniquely to a
codomain whose quotient is a Banach space
(Theorem~\ref{thm:operator-extension}), and from this we obtain the
operator-morphism version of the main theorem
(Theorem~\ref{thm:operator-main-theorem}).
\textbf{Here uniqueness does not follow automatically from the agreement of maps
of base spaces, but is obtained by way of contractivity and density.}
This is a difference from the non-ae version in
Section~\ref{chap:functoriality}.
The ae version of that section
(Subsection~\ref{sec:main-theorem-ae}) uses the same density.
We also see that, since the extension is determined only up to null functions,
it becomes natural to divide the codomain by null functions (we do not, however,
prove that ``the extension is undetermined unless one divides'';
Problem~\ref{prob:no-lift}).

Finally we investigate the relation between the two sections.
The functor from the category whose morphisms are maps of base spaces to the
category of operator morphisms becomes faithful once the ae-quotient is taken,
but is not full.
Using this, one finds that the \emph{uniqueness} part of the main theorem of
Section~\ref{chap:functoriality} follows from the theorem of
Section~\ref{chap:operator}, while \emph{existence} does not.
We also show that the morphisms of the category of operator morphisms see only
the quotient $\widetilde{L}$
(Proposition~\ref{prop:Q-full-faithful}).
Hence that category distinguishes neither the base space nor the realization as
a function space, and what the Daniell extension gives is the Banach completion
of a pre-AL space
(Proposition~\ref{prop:D-op-is-banach-completion}).
The subcategory to be taken as the codomain of the main theorem in this category
is the class of all objects whose quotient is a Banach space, not the class of
all Daniell closed objects.
The latter is not isomorphism-closed, and the former is its isomorphism closure
(Proposition~\ref{prop:BInt-characterization}).
\textbf{That is, only the side of Section~\ref{chap:functoriality} retains the
distinction between an abstract completion and a concrete realization as
categorical structure.}
This is the reason for keeping the two sections apart.

\medskip
\noindent
\textbf{The categories and functors of this paper.}
On the side where morphisms are maps of the base spaces there are two
reflections.
\[
  \begin{tikzcd}[column sep=4.5em, row sep=3.6em]
    \mathbf{DInt}
      \arrow[r, hook, "\iota_M"]
    & \mathbf{MInt}
      \arrow[r, hook, "\iota"]
    & \mathbf{Int}
      \arrow[l, bend right=42, "(-)^{\mathrm{M}}"']
      \arrow[ll, bend left=26, "\mathcal{D}"]
  \end{tikzcd}
\]
Both curved arrows are left adjoints:
$\mathcal{D}\dashv\iota_{D}$ (where $\iota_{D}=\iota\circ\iota_M$) and
$(-)^{\mathrm{M}}\dashv\iota$, the latter being $A\mapsto A^{\mathrm{M}}$
(Theorem~\ref{thm:mint-inclusion-has-left-adjoint}).
\textbf{These two reflections do not agree}, and that is the starting point of
Subsection~\ref{sec:why-complete}
(Remark~\ref{rem:D-I-not-left-adjoint}).

Dividing morphisms by null functions and then replacing them by operator
morphisms gives the following diagram.
\[
  \begin{tikzcd}[column sep=5em, row sep=3.4em]
    \mathbf{Int}_{\mathrm{ae}}
      \arrow[r, "\mathcal{D}_{\mathrm{ae}}"]
      \arrow[d, "\mathcal{E}_{\mathrm{ae}}"']
    & \mathbf{DInt}_{\mathrm{ae}}
      \arrow[d, "\mathcal{E}_{\mathrm{ae}}"]
    \\
    \mathbf{Int}_{\mathrm{op},\mathrm{ae}}
      \arrow[r, "\mathcal{D}_{\mathrm{op}}"']
      \arrow[d, "Q"']
    & \mathbf{BInt}_{\mathrm{op},\mathrm{ae}}
    \\
    \mathbf{AL}_0
  \end{tikzcd}
\]
The category $\mathbf{Int}_{\mathrm{ae}}$ in the upper left is obtained by
dividing the morphisms of $\mathbf{Int}$ by ae-equivalence
(Definition~\ref{def:int-ae}).
Both horizontal arrows are reflections
(Theorem~\ref{thm:main-ae}, Theorem~\ref{thm:operator-main-theorem}), and the
square commutes (Proposition~\ref{prop:E-commutes-with-D}).
The last functor $Q$ is faithful and full
(Proposition~\ref{prop:Q-full-faithful}), which is what the statement above
that the bottom row sees only the quotient $\widetilde{L}$ amounts to.

\medskip
The appendix treats the general theory of adjoint functors used in the proofs
and the idea of forgetting the base space.

\subsection{Relation to earlier work}\label{sec:prior-work}

\subsubsection{On the content of the preliminary section}

The content of Section~\ref{chap:preliminaries} follows \cite{Yamagami}.
The definition of integrability by series representations, the fact that $L^1$ is
a Riesz space, that $I^1$ is an integral, that the Daniell extension satisfies
the monotone convergence theorem, and the idempotence of the Daniell extension
are all in that book.
As that book itself states, these procedures follow \cite{BM}, and \cite{BM} in
turn points out that the construction of the complete Daniell space by this
method was already presented in \cite{M1989}.
The origin of the method itself is \cite{Mikusinski}.

That is, the lineage of the content stated in
Section~\ref{chap:preliminaries} is
\[
  \text{\cite{JM1974}},\ \text{\cite{Mikusinski}}
  \ \to\ \text{\cite{M1989}}
  \ \to\ \text{\cite{BM}}
  \ \to\ \text{\cite{Yamagami}}
  \ \to\ \text{this paper} ,
\]
and \textbf{this paper claims no novelty whatsoever for this section.}
Note that \cite{JM1974} has the same title as \cite{M1989} and is fifteen years
earlier, and may be the oldest source for the construction by series, but the
relation between the two could not be confirmed here.

Note also that the series representation adopted in this paper is the form of
\cite{Yamagami}, in which the dominating functions $\varphi_n$ are made
explicit, and differs on its face from the form without dominating functions
adopted in \cite{BM}.
That the two are equivalent is shown in
Lemma~\ref{lem:series-rep-equiv}.

\subsubsection{On arguments with the base space fixed}

The behaviour of the Daniell extension with the base space $X$ fixed, namely
that the extension is idempotent, that it is monotone with respect to inclusion
of extensions, and that it is the ``smallest complete extension'' (in the
terminology of this paper the smallest Daniell closed extension; see
Remark~\ref{rem:why-daniell-closed}), is stated as a corollary in
\cite{Yamagami}.
\cite{BM} makes remarks to the same effect.
There is also, under the name of the ``closure'' of a Daniell space,
\cite{CWS} and \cite{CFW}, which develop this kind of construction
systematically.

\textbf{Hence this paper claims nothing about minimality with the base space
fixed.}
What this paper treats is the part concerning what can be said when the base
space varies.

\subsubsection{On morphisms and categories}

A notion corresponding to a morphism between integration systems in fact also
appears in \cite{Yamagami}.
In the context of discussing iterated integrals, that book introduces linear
maps $\Lambda$ from a Riesz space $L$ on a set $S$ to a Riesz space $M$ on a set
$T$ that are positive and order continuous, and points out that these include
probability kernels, Markov operators and transition probabilities.
It also contains the fact that $\Lambda$ extends to $L^1$ for integrable
functions (although determined only up to null sets).

\textbf{What remains for this paper is the part that assembles these into a
category, and the part that decomposes Daniell closedness into two closure
conditions.}
Specifically, the following four points.

\begin{enumerate}
  \item Defining identities and composition, and verifying that the Daniell
    extension satisfies the functor laws.
  \item \textbf{Uniqueness of the factorization}. What earlier work supplies is
    the \emph{existence} of the extension, not that the morphism factoring
    through a given morphism is unique.
    Uniqueness is an assertion that acquires meaning only once a class of
    morphisms is fixed, and cannot even be posed as a question without
    introducing a category.

    Whether analytic content enters the uniqueness, however, depends on the
    choice of morphisms.
    When the maps of base spaces are taken as morphisms as they stand,
    uniqueness follows formally from the agreement of the maps of base spaces
    and has no content.
    When morphisms are divided by null functions
    (Subsection~\ref{sec:main-theorem-ae}) and when operator morphisms are
    adopted (Section~\ref{chap:operator}), an argument by way of the density of
    $L$ in $L^1$ becomes necessary.
  \item \textbf{Decomposing Daniell closedness into two closure conditions}
    (Theorem~\ref{thm:complete-intrinsic}).
    As long as $\mathbf{DInt}$ is defined by the condition $L^1=L$, the main
    theorem says no more than that the subcategory of fixed points of
    $\mathcal{D}$ has $\mathcal{D}$ as its reflector, and carries no new
    content.
    Through the equivalence of Daniell closedness with
    \begin{quote}
      satisfying the monotone convergence theorem and being closed under
      modification on null sets,
    \end{quote}
    the main theorem takes the form ``the reflector into the subcategory cut out
    by two conditions that occur independently in measure theory is given by the
    concrete construction of series representations''
    (Remark~\ref{rem:main-theorem-not-tautology}).
    This is not a matter of assembling a category but a part with analytic
    content.
    The line of argument ``it has content because it can be written without
    mentioning $L^1$'' does not, however, work.
    The reason is given in
    Remark~\ref{rem:not-about-mentioning-L1}.
  \item On the basis of the above, stating that $\mathbf{DInt}$ is a reflective
    subcategory.
\end{enumerate}

The four results treated in
Subsection~\ref{sec:why-complete}, namely the counterexample showing that the
Daniell extension is not left adjoint to
$\iota:\mathbf{MInt}\to\mathbf{Int}$, the construction of the reflection into
$\mathbf{MInt}$, the theorem determining that reflection by the
$\sigma$-algebra generated by $L$
(Theorem~\ref{thm:LM-equals-Lsigma}), and the decomposition of $L^1$
(Proposition~\ref{prop:l1-decomposition}), were obtained in this paper
independently of the above.
Taking the last two together, it follows that the gap between the two
reflections amounts exactly to the functions supported on null sets, and that for
the integration system of step functions this is the difference between Borel
and Lebesgue.
These, together with
Theorem~\ref{thm:complete-intrinsic}, are the parts of this paper with technical
content.
Both lie, however, in the neighbourhood of classical Daniell theory, and the
claim is only that \textbf{to the author's knowledge they have not been stated
explicitly}.

\subsubsection{Neighbouring work}

There is a body of work attempting to treat integration theory categorically.
\cite{Hartig} restates the Riesz representation theorem on compact spaces as a
natural equivalence between the functor giving the Banach dual of the space of
continuous functions and the functor giving the space of signed Baire measures.

In particular \cite{BorgerHabil}, and \cite{Borger2010}, which treats its
subject in a general category, construct the integral itself by a universal
property, starting from Boolean algebras.
For a Boolean algebra $\mathcal{B}$, the assignment of indicator functions is the
universal charge into the space of simple functions, and any charge factors
through it uniquely; the way it factors is what corresponds to ``the integral of
simple functions''.
The functor taking the algebra of simple functions is left adjoint to the functor
assigning to a ring the Boolean algebra of its idempotents.

As far as we have read the text of the published version \cite{Borger2010}, this
is a different theorem from the present one.
The starting point differs (a Boolean algebra, versus a pair consisting of a
Riesz space and an integral), the content of the universal property differs
(constructing an integral, versus extending an integration system already
given), and the destination differs as well (the space of simple functions,
versus the space $L^1$ of integrable functions).
The universal property in this paper does not produce an integral; it is a
reflection sending a given integration system to a Daniell closed one.

Moreover, the boundary between the two is made explicit within that paper.
What $LB$ represents is the space of simple functions, and immediately after its
Theorem 2.1 that paper states that, given suitable completeness conditions, the
integral can be extended from simple functions to more general measurable
functions, without carrying out that extension.
Concerning the monotone convergence theorem, it goes no further than remarking in
the introduction that convergence theorems such as monotone convergence and
dominated convergence ought to mean continuity in some topology.
\textbf{What this paper treats is precisely that side, the extension.}
Neither the Daniell integral, nor Riesz spaces, nor $L^1$ appears in that paper,
and there is no assertion about reflective subcategories or completions.
Its MSC classification, too, does not include 28C05, the subject of this paper.

We have not, however, obtained the text of \cite{BorgerHabil} (219 pages).
\cite{Borger2010} states that a preprint gives, on the basis of Linton's method,
a treatment in a special category similar to that of the
Habilitationsschrift, from which one can read off that that work treats the same
subject as above.
Even so, we record that \textbf{the possibility that results of the same kind are
contained there cannot be entirely ruled out}.

Work treating measure theory and probability from the categorical side has also
been active in recent years.
\cite{FritzLorenzin} regards Boolean $\sigma$-algebras and monotone
$\sigma$-complete C*-algebras as abstract measurable spaces and organizes
functors, adjunctions and equivalences with the category of measurable spaces.
What that paper treats, however, are measurable spaces, not objects carrying a
measure or an integral.
\cite{DahlqvistKozen} constructs a symmetric monoidal closed category of ordered
Banach spaces and regular operators and uses it for the semantics of
probabilistic programs.
\cite{EhrhardGeoffroy}, motivated by the semantics of probabilistic programming
languages, develops a theory of integration on cones modelled on the Pettis
integral and shows that the category whose morphisms are the integral-preserving
linear maps is a model of linear logic.
Its objects and morphisms differ from those of this paper, but it makes the same
choice in taking preservation of the integral as the condition on morphisms.
This reference was pointed out to the author by T.~Fritz; the author has read
only its abstract and has made no detailed comparison with its contents.

Also, the normed space obtained in this paper by dividing an object of
$\mathbf{Int}$ by the null functions is the pre-completion form of a space whose
norm is additive on the positive cone, that is, of an abstract Lebesgue space
(AL-space).
The operation of building this completion from a Riesz space and a positive
linear functional is itself standard and can be found in, for instance,
\cite{AbramovichAliprantis}.

Indeed, assertions of the form ``the completion is a reflector'' are a
well-established type in the theory of ordered vector spaces.
For example \cite{BMO} treats the Dedekind completion functorially, and
\cite{AvilesBilokopytov} states it explicitly as a reflector.
\textbf{What these treat, however, is the Dedekind completion, a completion
determined by the order structure alone, and it is a different completion from
the Daniell extension of this paper, which depends on the integral $I$.}
For the integration system of step functions, the Dedekind completion widens the
space in the direction of the order, whereas what the Daniell extension gives is
the space of Lebesgue integrable functions.

On the measure-theoretic side, all the ingredients corresponding to the
reflection of this paper are classically in place.
\cite{FremlinMT2} constructs the completion of a measure space in item 212C and
states in 212D that it is complete, that it extends the original measure, and
that it coincides with the original space if and only if that measure was
already complete.
That this extension is unique on the completed domain is 212E(c), and that the
integrals agree is 212F(b).
As for the class of morphisms, the definition of an inverse-measure-preserving
function is 234A, that the completion preserves this class is 234B(a), that the
class is closed under composition is 234B(g), and the transformation of
integrals along such a map is 235G.

\textbf{That book does not, however, assemble these into a functor, an
adjunction or a reflection.}
Indeed, none of the words ``functor'', ``adjoint'', ``reflect'' or ``category''
occurs anywhere in the relevant Chapters 21 and 23.
That is, the position taken in
Subsection~\ref{sec:our-position} --- \textbf{the ingredients are classical and
the assembly is this paper's contribution} --- is confirmed from the
measure-theoretic side as well.

Note that whether a completion is a reflector depends strongly on the choice of
the class of morphisms.
\cite{BMO} points out that the usual Dedekind completion is not functorial as it
stands and adds structure, and the reflection theorem of
\cite{AvilesBilokopytov} is likewise stated in a category of morphisms
preserving existing suprema.
It is for the same reason that this paper prepares two choices of morphisms and
discusses their relation in
Subsection~\ref{sec:two-chapters}.
What the main theorem of this paper states is not that \emph{an abstract
completion exists}, but that \textbf{the concrete construction of the Daniell
extension realizes that completion as a function space on the same base set, and
does so including the extension of morphisms}.
This distinction is meaningful in the case where Daniell closedness decomposes
into two conditions that occur independently in measure theory, the monotone
convergence theorem and modification on null sets
(Theorem~\ref{thm:complete-intrinsic}).
Under that decomposition the main theorem takes the form ``the reflector into the
subcategory cut out by these two conditions is given by the concrete
construction of series representations''
(Remark~\ref{rem:main-theorem-not-tautology}).
The formulation of Section~\ref{chap:functoriality} is given so that this point
is not lost.

\subsubsection{The position of this paper}\label{sec:our-position}

On the basis of the above, let us state the position of this paper honestly.

The results of Subsection~\ref{sec:why-complete} split the Daniell extension
into the following two stages
(Remark~\ref{rem:D-factors}).
\[
  L
  \quad\subset\quad
  L^{\mathrm{M}}=L^{\sigma}
  \quad\subset\quad
  L^1=L^{\mathrm{M}}+\mathcal{Z}_I
\]
The first stage is the operation of taking the integrable functions measurable
with respect to the $\sigma$-algebra generated by $L$, and the second is the
operation of completing the measure.
\textbf{Both are classical operations, and moreover it is this paper's own
theorems that show them to be so.}

Hence this paper is not ``a new completion theorem distinct from an abstract
completion''. In a phrase, it is

\begin{quote}
  \textbf{the Daniell counterpart of the completion of a measure space,
  organized in the language of categories}.
\end{quote}

\noindent
The four contributions listed up to
Subsection~\ref{sec:prior-work} are the breakdown of that organization.
The decomposition of Daniell closedness
(Theorem~\ref{thm:complete-intrinsic}) supplies the division into the first and
second stages; the determination of the second reflection
(Theorem~\ref{thm:LM-equals-Lsigma},
Proposition~\ref{prop:l1-decomposition}) identifies the destination of each
stage; the assembly into a category restates the composite as a reflector; and
the observations on the operator category show how far one can abstract before
the realization becomes invisible.

It is more accurate to put this position in the foreground than to hide it, and
it also makes it easier for the reader to see what is classical and what is not.

\subsubsection{On the scope of the literature survey}

In writing this paper we searched, via zbMATH Open, the intersection of
MSC 28C05 (integration theory via linear functionals) with 18-XX (category
theory).
Within that range we found no work identifying the Daniell extension as a
reflection into the Daniell closed integration systems.
The works obtained were \cite{Hartig} and \cite{BorgerHabil}, listed above, and
work on section spaces of Banach lattices, whose subject differs from this
paper's.

We also searched for the lineage described above, and confirmed from the review
of \cite{M1989} that that paper is the presentation of the construction by
series itself and contains no minimality, universality or categorical content.
We further searched for work treating abstract Lebesgue spaces and completions of
Riesz spaces as adjunctions or reflections; everything obtained concerned the
Dedekind completion, and we found nothing taking the Daniell extension as its
object.

For \cite{CWS} we consulted the book itself. What we checked is
Chapter 1, Section 4 (closed Daniell spaces), and the outcome is recorded in
Remark~\ref{rem:extended-real-collapse}.

\textbf{We have not, however, been able to carry out a more exhaustive survey.}
We state explicitly what has not been checked.

\begin{itemize}
  \item Of \cite{CWS}, \textbf{Chapter 4 (extensions of a Daniell space and the
    integral), Chapter 5 Sections 3 and 4 (measurability and integrability),
    and the historical remarks at the end of the volume have not been
    checked.}
    As far as one can judge from the table of contents and the index,
    Chapter 4 treats families of extensions and their order structure, which
    places it close to the content of
    Subsection~\ref{sec:why-complete}.
    In particular its Theorem 4.1.10 is, according to the index heading, a main
    theorem about the operation of adjoining exceptional functions, and may be
    an assertion of the same kind as
    Proposition~\ref{prop:l1-decomposition}.
    Likewise Chapter 5 Section 4 treats the relation between measurability and
    integrability and is the place to be compared with
    Theorem~\ref{thm:LM-equals-Lsigma}.
  \item \cite{CFW} is held by few libraries in Japan and we have not been able
    to consult its text.
  \item We have been unable to obtain the texts of
    \cite{BorgerHabil}, \cite{M1989} and \cite{JM1974}, and our grasp of their
    contents rests on the zbMATH reviews or is unconfirmed.
\end{itemize}

\noindent
Accordingly, \textbf{for the results of
Subsection~\ref{sec:why-complete}, and in particular for
Proposition~\ref{prop:l1-decomposition} and
Theorem~\ref{thm:LM-equals-Lsigma}, the comparison with earlier work is not
complete.}
If all or part of the results of this paper are already known, the author would
be grateful to be informed.
